% 独立审稿草稿；与论文译文分开，避免将审稿建议误写成已完成研究。
% 编译：tectonic Huang_FEHU_Midterm_Review_zh.tex（同目录须保留 fehu_paper_style.sty）。
\documentclass[12pt,a4paper]{article}
\usepackage{fehu_paper_style}
\hypersetup{pdftitle={FEHU 中期论文审稿草稿：证据、创新性与后续研究},pdfauthor={研究审读草稿},pdflang={zh-CN}}
\renewcommand{\contentsname}{目录}
\renewcommand{\figurename}{图}
\renewcommand{\tablename}{表}
\renewcommand{\refname}{参考资料}
\titleformat{\section}{\normalfont\sffamily\fontsize{14}{18}\selectfont\bfseries}{\thesection}{0.8em}{}
\titleformat{\subsection}{\normalfont\sffamily\fontsize{12}{16}\selectfont\bfseries}{\thesubsection}{0.8em}{}
\begin{document}
\begin{center}
{\fontsize{20}{26}\selectfont\bfseries 中期论文审稿草稿\par}
\vspace{8pt}
{\large FEHU：证据核验、研究深度、创新性与后续方案\par}
\vspace{8pt}
审读对象：Peiyang Huang 的 Midterm Paper\\
审读日期：2026 年 9 月 12 日
\end{center}
\noindent 本稿供作者自审和与导师讨论使用。结论依据提供的论文、仓库保留产物及所列一手资料，不代表学院或答辩委员会的正式判定。本文按\textbf{硕士中期研究}评价，不能直接套用博士资格审查标准。

\section{总体意见：具备中期答辩基础，但主张需要与证据对齐}
\textbf{建议结论：研究工作量和技术深度足以形成硕士中期答辩材料；以“完成了可验证的系统研究，并明确剩余科学问题”的定位进入答辩是合理的。当前证据不足以支持“增强稳定改善长尾识别”“全部组件必不可少”或“达到整体最先进水平”等更强主张。}如果作为最终学位论文或独立方法论文评审，我会要求重要修改，核心是补齐贡献归因、对照公平性和稳定性验证。

这并不是因为使用现成模型就缺少研究价值。该工作把多主体新闻、多标签长尾、部分监督和方向判断结合起来，包含实际训练、推理系统、官方评测和可解释的失败分析，明显超出简单调用 API 的演示。其短板是若干组成技术已经成熟，最具辨识度的组合机制尚缺少完整消融，而增强路线与最终融合收益之间的联系没有闭合。

原稿的一项优点是已经坦诚承认大部分限制。审稿重点应是把这些限制转化为少量可执行实验，而不是用更强措辞包装现有分数。若作者能够解释算法的真实行为、说明负面结果，并展示下述控制实验方案，我倾向于支持其作为中期进展通过；实际通过仍取决于口头答辩、课程理解和学院审查。

\subsection{与学院要求的对应}
早稻田 IPS 的硕士学习成果要求包含：适当设定课题、提出解决方法、客观验证其健全性、新颖性和有效性，以及准确交流研究观点。\href{https://www.waseda.jp/fsci/gips/about/policy/}{学院教育理念}为这些评价维度提供依据。2026 年指导材料说明，中期发表旨在检查硕士论文进展和课程理解程度，由含导师在内的三名教师审查；这不是一项“必须超过排行榜第一名”的公开要求。\href{https://www.waseda.jp/fsci/gips/assets/uploads/2026/02/Guidance-for-Master-course-students-Apr.2026_J-1.pdf}{2026 年硕士指导材料，第 49--51 页}。

若作者属于 2027 年 3 月毕业批次，学院另要求提交\textbf{一页 A4 的指定格式中期概要书}，截止时间为 2026 年 9 月 30 日 17:00（日本时间）。本次中英文全文用于研究内容审核，不能替代该概要书。应按\href{https://www.waseda.jp/fsci/gips/other/2026/07/09/35430/}{本轮学院通知}处理正式表格和签章；本次没有代为提交任何材料。

\begin{table}[htbp]
\centering\small
\caption{按硕士中期研究要求形成的定性判断；不是学院评分表。}
\begin{tabularx}{\textwidth}{@{}P{2.5cm}P{2.3cm}Y@{}}
\toprule
维度 & 判断 & 判断依据及缺口 \\
\midrule
问题价值 & 满足 & 行为主体、价值观和方向的联合判断具有明确任务边界。 \\
技术深度 & 基本满足 & 掩码监督、LoRA、交叉校准、集合决策与方向验证需要实质性实现和理解。 \\
已完成工作量 & 满足 & 有三种子训练、多个系统、分阶段结果和官方提交；运行数量本身不等于创新。 \\
方法创新性 & 增量创新 & 有任务化组合设计，但缺少对其不可替代性的直接证明。 \\
实验可信度 & 核心数值可核 & 开发结果可重算；测试成绩只能与所提供官方表格核对。 \\
因果归因 & 尚不充分 & 无法把完整融合收益分配给增强、编码器或每个评审成员。 \\
终稿完成路径 & 可行 & 优先做固定输出消融，再推进质量审核和 Task~2；避免同时扩张全部路线。 \\
\bottomrule
\end{tabularx}
\end{table}

\section{真实性与数值证据：哪些已经核实}
本次使用当前仓库的官方评估器对保留预测重新评分，并从逐次运行 CSV 重新计算统计量；没有重新训练模型或重新调用外部生成服务。因此，下面的“复核”指\textbf{现存数据与预测的数值复核}，并不等于重新完成整个实验过程。

\begin{table}[htbp]
\centering\small
\caption{本次核心复核结果。原稿表号在中英文版中保持一致。}
\begin{tabularx}{\textwidth}{@{}P{3cm}YP{3.3cm}@{}}
\toprule
对象 & 核验结果 & 证据边界 \\
\midrule
数据规模（表 4） & 训练 1,228 篇／3,381 个给定主体；开发 177／459；Task~1 测试 175／479；Task~2 测试 175／444。 & 给定主体与有正例主体不同。 \\
标签统计 & 训练 16,905 个带方向元组；高频门控包含 38 类；三类训练支持度为零。 & 元组数不是独立文章数。 \\
E7 各阶段（表 9） & G0/G1/G2/G3/G4/G5 的 Micro-F1、Macro-F1、DRR 和输出计数与原稿一致。 & 固定的七月预测文件。 \\
三种子结果（表 7） & A/B/C/Fp/Gp/N1/N3 共 21 行运行记录的均值和样本标准差与原稿舍入值一致。 & 从保留 CSV 复算，未重训。 \\
主结果置信区间 & 重跑 10,000 次文章级配对自助法，G4 减 G0 为 0.0321，95\% 区间 [0.0209, 0.0433]。 & 条件于现有输出与已选方法。 \\
官方成绩（表 10） & 所提供任务概述的 Task-1b 表确认 WAIPS-03 为 0.5808／0.3952，组织者为 0.6284／0.4402。 & 无本地测试金标准，未独立重算测试成绩。 \\
提交覆盖 & WAIPS-02：175 篇、479 个非空主体、2,260 个带方向元组；WAIPS-03：175 篇、470 个非空主体、2,764 个元组。 & 非空覆盖不同于给定主体覆盖。 \\
增强质量说明 & 第二轮清单确认 900 条记录中，450 条完整监督、450 条部分监督，只有 36 条带显式已知负例。 & 结构验证不证明语义正确。 \\
\bottomrule
\end{tabularx}
\end{table}

在已核对项目中，未发现核心结果的计算不一致，所列分数具有实际文件支撑。仍不能据此认证每次 API 请求的服务端模型版本、全部历史数据处理过程、所有合成标签的正确性，或全部检查点的可复现性。作者应把“数值可复算”“运行可重现”“语义可信”分别说明，避免把它们合并为一句“实验真实可靠”。

官方结果的排名表述必须保持准确：WAIPS-03 的 Task-1b Macro-F1 在参赛提交中最高，但只有两个参赛团队、四个参赛包，且组织者基线更高。其优势并不覆盖全部指标。相对组织者的 Micro-F1 差距是 0.0476，即 4.76 个百分点，不能用“取得第一”掩盖。

\subsection{本次补充分析：方向检查的净收益较小}
对 G3 与 G4 的标签变化逐项核对，46 次方向翻转中，14 次由错误变正确，7 次由正确变错误，5 次前后均受金标准支持，20 次前后均不受支持。因此净增加 7 个正确元组；“翻转了 46 个方向”不能写成“修复了 46 个错误”。

本次对 G4 与 G3 单独重跑相同的 10,000 次文章级配对自助法，Micro-F1 差值为 0.0027，95\% 区间为 \textbf{[-0.0008, 0.0062]}，包含零。这是本次审稿新增的固定文件分析，未被写回忠实译稿。它支持把方向问答描述为本次输出上的小幅改进，尚不足以声称方向组件具有稳定显著收益。该区间仍没有计入反复选模型或外部生成变化。

\section{主要问题及修改要求}
\subsection{优先级 P0：增强与最终融合之间缺少直接证据}
位置：摘要、第 1.3、4.2、4.6、5.6、5.7 和 5.11 节。18 个增强目标全部满足训练计数不超过 130，而删除和添加只作用于计数大于 130 的标签。于是 G0、G1、G3 的目标 Micro-F1 同为 0.3145；G4 的 0.3264 来自方向变化。这里没有数学矛盾，但存在\textbf{读者预期与证据之间的缺口}：标题突出增强，最强结果却是对主要由其他标签构成的系统输出进行修订。

最低修改要求是把三个问题分开：增强是否改善目标排序；编码器是否给融合提供了有用信息；增强编码器是否优于仅官方数据编码器。完整 FGHF 至少比较无编码器项、仅官方数据编码器、增强编码器三种来源。固定主预测、评审成员、文章折、门控及调用预算，并为每种条件重新拟合校准器。不能简单将编码器置零却保留原校准器，然后把性能变化都归因于信息缺失。

若无编码器版本相当，应接受“有用的融合不依赖增强编码器”这一结果，简化系统，并把增强作为重要的失败分析或独立研究问题。当前题目可保留用于中期汇报；若终稿仍不能证明增强贡献，建议调整题目重心，避免让增强与已经证实的融合贡献处于同等承诺位置。

\subsection{优先级 P0：开发集选择偏差没有被两折校准消除}
位置：第 4.6、5.1、5.7、5.12 节。目标标签、提示、早停、门控、分数规则和最终配置都受同一开发集影响。两折交叉拟合只隔离当前校准步骤，不会撤销此前的选择。现有 G4--G0 自助法区间说明固定文件之间有差异，不能认证整套研究流程的泛化。

已有官方测试成绩是独立评估的有价值观察，但不能用于继续调参后仍声称测试完全未触及。现在从已经反复查看的 177 篇文章中再划一部分，不能使其追溯性地变成未见测试集。后续可在原开发集上报告探索性分析，并对新收集、锁定且未用于设计的评估数据进行确认；若采用外层交叉验证，所有可执行的数据选择和校准应进入内层，仍需披露人工设计曾使用整个开发集的历史。

\subsection{优先级 P1：增强质量与监督目标要分别验证}
位置：第 4.2、4.3、5.6 节。第一轮自动语义筛查、第二轮结构检查、六条人工试点以及 90 条“困难负例”生成类别，不能互相替代。尤其不能把 1,240 条记录统称为“人工验证的增强集”。原稿对此的限定是正确的，应保留。

Fp 的两个种子在全局阈值下没有目标真阳性；N1 虽有 0.9381 的平均召回，精确率仅 0.0179。这说明必须同时报告排序、工作点和假阳性成本。对长尾标签建议补充 PR-AUC／平均精确率，并列出各标签正例数；ROC-AUC 不足以体现稀有正例的实际检索代价。阈值只能在训练／校准部分选择。

掩码损失按受监督条目总数归一化，因此每条完整监督记录与每条部分监督记录的总贡献可能不同。负例下采样还同时改变有效先验、正例权重及归一化分母。论文应说明这些耦合，不能仅以“原先负例太多”解释全部变化。建议先按每目标十条开展分层审核，至少两位审核者独立判断主体、价值观、方向和未知标签；这约 180 条的审核适合发现机制性错误，但不能被包装为每类精确的总体质量估计。

\subsection{优先级 P1：融合规则合理，但必要性与校准可靠性未证实}
位置：第 4.6--4.7、5.7、5.11 节。频次门控的动机合理，但相对于只有删除分支的 G1，G5 去除门控后 Micro-F1 更高、Macro-F1 更低，说明它是目标权衡，不能宣称普遍优于无门控。G2、G5 只有删除分支，不能代替完整 G4 的组件消融。

还需重点处理三项机制：第一，添加校准以“原始主集合没有该价值观”为候选条件，而应用以“删除后没有”为条件，导致候选分布不匹配；第二，每主体至少保留一个标签可能强迫产生假阳性，任务允许空集；第三，最多两条新增元组和两个方向的处理是启发式规则，需要报告影响范围。原稿准确揭示了这些实现细节，这比将算法美化成理想版本更可信。

后续优先验证完整系统有／无门控、等预算简单投票及校准器自身质量。可报告 Brier 分数、带样本量的可靠性图，并以严格内层选择的简单阈值作为对照。$F_1/2$ 在这里是启发式阈值，不是已证明的全系统最优阈值；其理论来源的适用条件应继续保留。\href{https://arxiv.org/abs/1402.1892}{Lipton 等，2014}。

\subsection{优先级 P1：强基线、公平预算与复现控制仍不足}
位置：第 5.4、5.8、5.10 节。组织者 8B 模型与 DeBERTa 在骨干、上下文、训练粒度和目标上同时不同。其更强成绩构成需要认真面对的基线，但不能直接证明其中某一个因素有效。应在同一模型下比较文章监督与文章加子事件监督，并记录实际截断比例；512 token 输入能否包含关键主体和证据，应由数据核验，而非仅靠猜测。

WAIPS-02 与 E7 的测试逻辑请求数分别为 700 和 1,432，不能仅比较 F1 就宣布更优方法或更高成本效益。应记录输入／输出 token、重试、总延迟，并加入相同预算下的简单集成。编码器还应统一对全部给定主体输出分数，包括无正例主体，以免比较受 456 与 459 的输入覆盖差异影响。

七月至八月同时改变多项条件，现有下降不能被命名为纯粹“API 漂移”。未来每次运行应保存数据、提示、模型参数和概率文件哈希；缓存键须反映这些条件；请求失败应明确记录并阻断不完整的新比较。历史失败记录应如实保留，不能静默补跑后继续沿用旧分数。

\subsection{优先级 P1：行为主体归属仍是未被充分测量的核心}
位置：第 3.2、5.9 和 6.4 节。两个案例能解释方向机制，不能估计总体归属能力。预测标签正确，也不代表解释使用了正确行为主体或证据。要让“actor-conditioned”成为实证贡献，需要一个独立审核的错误样本，区分主体错误、价值观错误、方向错误、证据不足以及合理的多重解释。

可以对有效样本实施主体替换、移除支持句、保留无关句等干预，但要人工确认干预后的标签确实应变化或保持不变。模型自己生成的解释不能再由它自己认可后作为事实真值。应把最终标签指标与证据支持率分别报告。

\section{创新性评价：有任务化设计贡献，尚非新的通用算法}
基于对主要一手文献的针对性核对，本文的合理定位是\textbf{面向特定任务的增量方法与实证研究}。不能因为合成数据、LoRA、掩码或投票早已有之，就否认全部研究价值；但也不能把组合了这些技术等同于发现了新的学习原理。本次核对不是穷尽性查新，不支持“领域首次”的断言。

\begin{table}[htbp]
\centering\small
\caption{创新性应落在何处，以及还需要什么证据。}
\begin{tabularx}{\textwidth}{@{}P{2.5cm}YY@{}}
\toprule
研究要素 & 与已有研究的关系 & 可以争取的贡献 \\
\midrule
价值观任务 & ValueEval 已研究层级价值观；NEVU 明确研究主体、多粒度与方向。 & 不能把任务本身作为作者首创；贡献是该任务中的具体决策设计。 \\
生成数据加本地模型 & EAVIT 已结合解释式训练、生成数据、本地价值观检测器和外部模型。 & DRA 的定向选择、监督语义及其失败机制需要独立证据。 \\
掩码与 LoRA & 部分标签掩码及参数高效适配均有成熟先例。 & 准确使用并分析适用条件，属于扎实实现，不应单独称新算法。 \\
频次门控异构融合 & 投票、保序校准和频次分组都已有基础。 & 删除／添加证据不对称、频次保护和方向触发的任务化组合，需超越简单对照。 \\
负面结果分析 & 负结果不是创新性的自动保证。 & 若受控证明数据构成与监督方式如何交互，可形成可迁移的研究认识。 \\
\bottomrule
\end{tabularx}
\end{table}

相关一手依据为 \href{https://aclanthology.org/2023.semeval-1.313/}{ValueEval（2023）}、\href{https://arxiv.org/abs/2603.17838}{NEVU（2026）} 和 \href{https://www.ijcai.org/proceedings/2025/934}{EAVIT（IJCAI 2025）}。尤其 EAVIT 已用本地检测器区分确定相关、确定无关和不确定价值观，再把不确定项交给外部模型；与本文的关系应在算法层面比较，而不只是列出引用。

建议用一个集中问题组织终稿：“在行为主体条件化的长尾多标签任务中，什么证据适合增加标签，什么证据适合保留或删除标签，何时值得额外检查方向？”如果同等预算下简单投票已经相当，复杂融合就应简化；如果选择性决策持续获益，就用分组错误和消融解释原因。这样的研究问题与现有证据接近，也比笼统宣称“增强和集成提高识别性能”更容易答辩。

\section{Task 2：作为有明确假设的后续研究方向}
\subsection{为什么值得加入}
Task~2 以文章--子事件--行为主体为单位，提供官方子事件，可以更直接地观察局部行动、主体和方向之间的联系。它能补上当前文章级指标无法充分检查归属的缺口，并检验 Task~1 的自主分段是否产生了有用中间表示。官方定义见\href{https://github.com/WangYao-GoGoGo/ntcir19_fehu}{FEHU 官方资源库}。本文现有结果仍然仅属于 Task~1；新增 Task~2 应明确标为计划。

\textbf{值得研究的假设是：子事件级监督或可核查的局部证据，能否改善主体归属与方向判断，并在不增加测试时特权信息的条件下迁移到文章级识别？}仅把原提示换成子事件输入并取得一组分数，通常只是任务扩展，创新性有限。更有价值的是说明局部证据何时有效、何时不足，以及系统如何处理这一差别。

\subsection{本次数据诊断说明不能采用简单硬并集}
本次将官方训练／开发金标准中的子事件标签投影为文章--行为主体--价值观--方向集合，再与文章级金标准比较，得到表~\ref{tab:scopeaudit}。这是结构诊断，没有训练 Task~2 模型，也不是 Task~2 性能结果。

\begin{table}[htbp]
\centering\small
\caption{文章金标准与子事件金标准并集的关系，均按唯一元组计数。}
\label{tab:scopeaudit}
\begin{tabularx}{\textwidth}{@{}Yrrrr@{}}
\toprule
划分 & 文章级金标准 & 子事件投影并集 & 二者交集 & 仅文章级存在 \\
\midrule
训练 & 16,905 & 12,788 & 12,788 & 4,117 \\
开发 & 2,390 & 1,799 & 1,799 & 591 \\
\bottomrule
\end{tabularx}
\end{table}

开发集直接使用这一金标准并集，文章级召回也只有 $1799/2390=75.27\%$；另外 24.73\% 的文章元组未在局部金标准并集中出现。数据现象并不自动说明这些标签为什么缺失，可能涉及标注范围或跨句整合，需人工检查。它已经足以否定“文章金标准必然等于子事件金标准并集”的实现假设。

因此，后续可保留局部证据分支，同时允许全文分支补充文章级判断；跨粒度一致性宜作为需要验证的软约束或诊断信号。不能用硬并集强制删除所有缺乏单一子事件支持的文章标签，也不能从该金标准诊断直接推断实际模型的性能上限。

\Needspace{15\baselineskip}
\subsection{最小实验设计及信息边界}
\begin{table}[H]
\centering\small
\caption{建议的 Task~2 与跨粒度最小实验，不预设任何条件必胜。}
\begin{tabularx}{\textwidth}{@{}P{1.2cm}YY@{}}
\toprule
条件 & 设计 & 能回答的问题 \\
\midrule
T0 & 固定同一个 8B 级骨干，只用文章级监督训练，在 Task~1 上评估。 & 建立较强且可控的文章基线。 \\
T1 & 同骨干加入官方训练子事件监督，同时报告 Task~1 和 Task~2。 & 额外监督粒度是否有迁移价值；需匹配训练 token／更新预算。 \\
T2 & 基于 T1 输出可核查的主体--行动--证据--方向记录，与普通标签输出作同预算比较。 & 显式证据是否改善归属，还是只增加文本和成本。 \\
T3 & 在 T2 基础上，按跨粒度分歧选择少量额外检查。 & 是否优于等预算随机检查；与全面检查相比，性能与成本如何权衡。 \\
\bottomrule
\end{tabularx}
\end{table}

优先完成 T0--T1，再决定是否继续 T2--T3。训练对照既要报告相同数据预算，也要说明额外子事件引入的重复上下文和额外训练量。所有同一文章的文章级、子事件级、主体级实例必须进入同一个数据折。不得用同一文章的 Task~1 标签训练，再把该文章的 Task~2 标签当作独立外部验证。

Task~1 测试时仅使用其给定文章和主体，自主分段必须由文章生成；Task~2 使用其官方给定子事件及标识符。两个官方测试集的文章不同，不能将 Task~1 与 Task~2 分数差解释为同样本的粒度效果。可在同一开发文章上做配对诊断，但使用官方子事件的 Task~1 条件必须明确标为额外信息诊断，不得与标准 Task~1 结果混排。

预先指定主要指标，例如 Task-2b Micro-F1，并共同报告金标准正例支持的 Macro-F1、DRR、预测覆盖、双方向样本、主体错误率和证据支持率。跨粒度迁移另报 Task-1b。通过文章级配对区间表达不确定性，记录实际 token 和请求成本。即使 F1 没有改善，若能可靠说明跨粒度失败边界或证据错误类型，也可以形成有价值的论文结果。

\Needspace{23\baselineskip}
\section{逐部分修改意见}
\begin{table}[H]
\centering\small
\caption{面向下一轮内容修订的清单。本次忠实译稿未擅自加入这些新主张。}
\begin{tabularx}{\textwidth}{@{}P{2.3cm}YY@{}}
\toprule
位置 & 当前评价 & 下一轮具体修改 \\
\midrule
标题与摘要 & 结果与限制已较诚实，但标题对增强的承诺偏强。 & 若完整消融无增强收益，调整标题重心；加入主指标、官方比较边界和主要缺口。 \\
引言／RQ & 研究问题清晰，但三条路线容易显得并列堆叠。 & 给出一条决策问题主线，并明确每条路线回答哪一个可证伪问题。 \\
相关工作 & 关键引用真实，组合技术的先例说明较准确。 & 增加与 EAVIT 的输入、候选选择、监督和调用预算对照；精简通用背景。 \\
方法／算法 & 数学与实现基本一致，边界说明充分。 & 显式说明排序同分规则、候选校准差异、未知标签贡献及空预测约束。 \\
实验设计 & 已区分开发、官方测试与后续运行。 & 用一张划分流程图交代所有选择位置；统一主体覆盖和基线预算。 \\
结果与案例 & 主分数可靠，但案例不能支撑归属能力结论。 & 增加 QA 转移矩阵、组件消融和错误样本审核；弱化仅靠翻转数的论据。 \\
结论与计划 & 已承认增强未闭合，方向正确。 & 把最少必要实验、成功标准与失败后的简化选择写成明确里程碑。 \\
附录／书目 & 路径有用，但不能代替来源版本。 & 为决定性文件补充版本或哈希；将尚处准备阶段的会议稿与正式出版区分。 \\
\bottomrule
\end{tabularx}
\end{table}

\section{建议的执行顺序与完成标准}
\begin{enumerate}
\item \textbf{答辩前：整理一个能独立核验的证据包。}固定当前主要结果文件和评估脚本，展示主结果、三种子表、完整系统缺失消融及 QA 净变化。准备一页“已证实／未证实”说明，并按学院要求制作正式概要书。不要为赶答辩并行启动多条新模型路线。
\item \textbf{随后 2--4 周：完成固定输出的核心消融。}先无编码器／官方编码器／增强编码器，再完整系统有／无门控和一个等预算简单投票对照。重新拟合各条件校准器；将原开发集结果如实标为探索性。交付物是少量条件清晰、可重跑的比较及其解释。
\item \textbf{随后 4--8 周：完成增强质量与强基线验证。}审核约 180 条合成样本；使用相同种子和可比预算比较原增强、审核子集、仅官方数据。建立同骨干文章与文章加子事件训练基线。若收益不稳定，应收缩增强主张，而不是继续无限生成数据。
\item \textbf{随后 8--12 周：有条件推进 Task~2。}优先 T0--T1 与主体归属审核，再按结果决定是否开展显式证据和选择性检查。以锁定的新评估数据或严格披露局限的外层评估形成最终证据。具体周期需依据 GPU、模型调用和人工审核资源调整。
\end{enumerate}

最小合格终稿不必让所有模块都显著获胜，但应对一个明确问题给出可信答案：例如，选择性融合确实有效，而增强不是必要来源；或经过审核的部分监督能稳定改善一个可用工作点。相反，持续增加实验名称，却无法分离变化原因，会削弱终稿的可信度。

\section{答辩中最可能被追问的问题}
\begin{description}[style=nextline]
\item[你的创新究竟是什么？] 回答任务化的选择性证据决策及其受控评估；承认投票、掩码、LoRA 和保序回归不是个人首创。
\item[为什么标题强调增强，但最终增益不一定来自增强？] 说明这是中期仍未闭合的假设，并展示三种编码器来源的完整系统消融计划；不能把模型用过增强数据当成因果证明。
\item[组织者 8B 模型为什么比你的系统强？] 给出实际差距，并说明骨干、上下文、目标和监督粒度同时不同；目前不能确定原因。
\item[两折校准是否已经排除过拟合？] 只排除了校准步骤中部分直接信息复用；人工选择历史和方法选择仍在，需要新评估或严格限定的外层实验。
\item[掩码是否有效？] 监督语义合理，但预测收益依赖数据构成且种子不稳；用 B/C 与 Fp/Gp 的结果回答，而非绝对肯定或否定。
\item[方向检查修复了多少错误？] 46 次翻转中 14 次改善、7 次损害，净增加 7 个正确元组；单独增益区间包含零。
\item[为什么不直接做 Task~2 或取子事件标签并集？] Task~2 可以强化局部归属研究，但文章与子事件标签不完全相同，且测试输入边界不同；开发集并集缺少 591 个文章元组。
\item[如果去掉编码器性能不降，研究是否失败？] 不是。它说明真正有效的信息来自其他组件，应简化系统，并把增强失败的条件解释清楚。
\item[你个人完成了什么？] 对照 DRA-FEHU、WAIPS-02 和 WAIPS-03/E7-G4 的设计、实现和分析回答；明确辅助定向提示与 Liu Dongyang 的完整 WAIPS-01 之间的边界，保留共同工作的署名。
\end{description}

\section{语言、格式及证据文件说明}
中英文稿保留相同的六个正文章节、一个附录、36 个小节、13 张表、6 幅图、1 个算法和 16 条参考文献。译文使用“行为主体”“一致／相悖”“部分标签监督”“保序回归”“频次门控”等一致术语；文献标题、模型标识符和路径保留可检索形式。陈列公式、交叉引用标识和表中实验值逐项对应。

全文采用 A4 单栏，四边页边距 25.4 mm；正文 12 pt，英文为 Times New Roman，中文为宋体，中文标题使用黑体。正文行距系数为 1.15，一级标题 14 pt，二级标题 12 pt；图表与图注 10.5 pt，路径密集附录 10 pt；数学使用与西文衬线字体协调的 NewTX 数学字体。表题位于表格上方，图注位于图下方，表格采用三线表形式，目录与交叉引用自动生成。这是一套通用全文论文格式，不宣称是学院指定概要书模板。

除排版外，只做了可核实的编辑性修正：统一参赛队名为 IIUoT；修复附录中失效的任务概述本地路径；将消融表中与实验问题相反的“negative scarcity”改为“excess negative supervision”；把 NTCIR-19 的书目状态写为会议准备稿，并将“published”类措辞规范为已报告结果。依据官方日程，2026 年 9 月属于论文草稿阶段，最终稿截止与会议在其后，不能仅凭本地 PDF 称其已经正式收入论文集。\href{https://research.nii.ac.jp/ntcir/ntcir-19/index.html}{NTCIR-19 官方日程}。没有更改实验分数、增补未经验证的实验结果，或将本审稿的新建议混入译文。

关键复核入口相对于 \path{ntcir19_fehu-master/}：
\begin{itemize}
\item 指标定义：\path{evaluation.py}；数据：\path{dataset/gold_labels/} 和 \path{dataset/hv_categories/}。
\item 三种子统计：\path{new_check/phase2/outputs/per_run_metrics.csv}、\path{new_check/phase3/outputs/per_run_metrics.csv}。
\item 融合预测与问答：\path{RH_test/E7_calibrated_jury_fusion/runs/20260725_004522_G4_dirqa/} 及同级各阶段目录。
\item 自助法：\path{RH_test/E7_calibrated_jury_fusion/scripts/bootstrap_e7.py}；本次使用 10,000 次重采样，随机种子 20260725。
\item 官方成绩依据：本目录的 \path{NTCIR_FEHU_TaskOverview (1).pdf} 第 4 页，以及 \path{WAIPS at the NTCIR-19 FEHU.pdf}。
\item 提交文件：\path{DS_final/submission/}、\path{E7G4Test/submission/}。
\end{itemize}
本次复算的两个区间文件保存在工作区的 \path{tmp/pdfs/midterm_review/} 下，分别为 \path{bootstrap_recomputed.json} 和 \path{qa_bootstrap_recomputed.json}；跨粒度表直接从上述金标准按文章、主体、价值观、方向去重构造。原始中期英文源稿和 PDF 同样保留在该临时审核目录，便于对照本次格式修订。
\end{document}
